\section{Divisor sheaves, Artin chains, jet supports, and packet descent}
\label{sec:divisor-sheaf-audit}

This section reconstructs one programme from the local notes as a sequence of
typed constructions.  The classical inputs are deliberately narrow.  Lee gives
the local factorization of a meromorphic function on a Riemann surface, the
associated divisor language, skyscraper sheaves, sheafification, and the
acyclicity of flasque sheaves
\cite[pp.~93--95; Chap.~5, especially Thm.~5.9; pp.~184--188, Thm.~6.34]{lee2024}.
Vakil gives a second formulation of sheafification by compatible germs and the
one-dimensional regular-local-ring/DVR equivalences
\cite[\S\S2.2.9 and 2.4.4--2.4.7; \S13.5]{vakil2024}.
All special constructions and all comparison maps below are proved here.  In
particular, neither classical source is being cited for the Boolean shadow or for
the chain--cube adjunction.

\subsection{The literal fixed pair at a zeta zero}

Retain the split-zero semiring $S=G(\C)=\C\sqcup\{\tau\}$ of
\cref{sec:gcue-audited-core}.  Its additive zero is $\tau$.  To avoid confusing
that element with the ordinary complex number zero, write
\[
 e:=0_{\C}\in\C\subset S,
 \qquad
 A:=\{\tau,e\}\subset S.
\]

\begin{proposition}[Fixed pair in literal coordinates]
\label{prop:ds-fixed-pair}
The operations inherited from $S$ make $A$ a unital semiring, with additive
zero $\tau$ and multiplicative identity $e$.  The map
\[
 \beta:A\longrightarrow\mathbb B=\{0,1\},
 \qquad \beta(\tau)=0,\quad \beta(e)=1,
\]
is a semiring isomorphism.  If $u=-1\in\C\subset S$, then
\[
 \operatorname{Fix}(u\cdot-)=A.
\]
Consequently, at every zero $\rho$ of $\zeta$,
\[
 \zeta(\rho)-1=-1,
 \qquad
 \operatorname{Fix}\bigl((\zeta(\rho)-1)\cdot-\bigr)=A.
\]
\end{proposition}

\begin{proof}
The four sums and products in $A$ are
\[
 \tau+\tau=\tau,\quad \tau+e=e+\tau=e,\quad e+e=e,
\]
and
\[
 \tau\tau=\tau e=e\tau=\tau,\qquad e^2=e.
\]
These are exactly the Boolean tables under $\beta$.  Multiplication by $u$ fixes
$\tau$ because $\tau$ is absorbing.  For a supported element $z\in\C\subset S$,
$uz=z$ is the ordinary equality $-z=z$, hence $2z=0$ and therefore $z=0_{\C}=e$.
This proves the fixed-set identity.  The final assertion is the direct substitution
$\zeta(\rho)=0$; it detects no further analytic property of $\rho$.
\end{proof}

The internal unit of $A$ is $e$, whereas the unit of $S$ is $1_{\C}$.  Thus the
inclusion $A\subset S$ is deliberately not being called a unit-preserving semiring
map.

\subsection{A product sheaf and the exact sheafification map}

Let $X$ be a Riemann surface and let
\[
 D=\sum_{x\in Z}m_x[x]
\]
be an effective divisor with closed discrete support $Z\subset X$ and
$m_x\in\mathbb Z_{\geq1}$.  For an open $U\subset X$, put
\[
 P_D(U):=\bigoplus_{x\in U\cap Z}A^{m_x},
 \qquad
 \mathcal V_D(U):=\prod_{x\in U\cap Z}A^{m_x},
 \qquad
 \mathcal A_D(U):=\prod_{x\in U\cap Z}A.
\]
The direct sum means finite support relative to the basepoint
$(\tau,\ldots,\tau)$; every restriction map forgets the coordinates outside the
smaller open set.

\begin{lemma}[Product sheaf on discrete support]
\label{lem:ds-product-sheaf}
For any family $(M_x)_{x\in Z}$ of sets with fixed algebraic structure,
\[
 \mathcal F_M(U):=\prod_{x\in U\cap Z}M_x
\]
is a sheaf with that structure.  Its stalk is $M_x$ at $x\in Z$ and the terminal
one-element object at $y\notin Z$.
\end{lemma}

\begin{proof}
For an open cover $U=\bigcup_iU_i$, compatible families on the $U_i$ specify
one and only one value at each $x\in U\cap Z$: choose any $i$ containing $x$,
and compatibility makes the choice irrelevant.  These values form the unique
global family.  If $x\in Z$, discreteness gives a neighbourhood $W$ with
$W\cap Z=\{x\}$, so the directed system defining the stalk is eventually
constant with value $M_x$.  If $y\notin Z$, closedness gives a neighbourhood
disjoint from $Z$, and the directed system is eventually terminal.
\end{proof}

\begin{theorem}[Finite support sheafifies to arbitrary support]
\label{thm:ds-sheafification}
The coordinate inclusion
\[
 \eta_D:P_D\longrightarrow\mathcal V_D
\]
exhibits $\mathcal V_D$ as the sheafification of $P_D$ in sheaves of
$A$-semimodules.
\end{theorem}

\begin{proof}
The target is a sheaf by \cref{lem:ds-product-sheaf}.  Let $\mathcal G$ be a
sheaf of $A$-semimodules and let $\varphi:P_D\to\mathcal G$ be a presheaf
morphism.  Fix $s=(s_x)_{x\in U\cap Z}\in\mathcal V_D(U)$.  For every
$x\in U\cap Z$, choose $U_x\subset U$ with $U_x\cap Z=\{x\}$, and put
$U_0=U\setminus Z$.  The family
\[
 \{U_0\}\cup\{U_x:x\in U\cap Z\}
\]
covers $U$.  On $U_x$, the section $s$ is the finite-support section $s_x$;
on $U_0$ it is zero.  If $x\ne y$, then
$(U_x\cap U_y)\cap Z=\varnothing$, and the same is true of
$U_x\cap U_0$.  Hence
\[
 \varphi_{U_x}(s_x),\qquad \varphi_{U_0}(0)
\]
agree on every overlap.  They glue uniquely in $\mathcal G$; call the result
$\widetilde\varphi_U(s)$.  The construction commutes with restriction and with
the $A$-semimodule operations because both statements hold on the displayed
cover and sections of a sheaf are determined locally.  Thus
$\widetilde\varphi:\mathcal V_D\to\mathcal G$ is a sheaf morphism.

If $s$ already lies in $P_D(U)$, the glued local images are precisely the
restrictions of $\varphi_U(s)$, so
$\widetilde\varphi\eta_D=\varphi$.  If another extension has this property, it
agrees with $\widetilde\varphi$ on every $U_x$ and on $U_0$, hence everywhere.
This is the universal property of sheafification.  It also agrees with the
compatible-germ construction in \cite[\S2.4.4--2.4.7]{vakil2024}.
\end{proof}

\subsection{Boolean coordinates, reduction, and lost information}

For $m\geq1$, let $[m]=\{1,\ldots,m\}$ and $B_m=\mathcal P([m])$, ordered by
inclusion and equipped with union as join.  Define
\[
 \Theta_m:A^m\longrightarrow B_m,
 \qquad
 \Theta_m(a)=\{r:a_r=e\}.
\]

\begin{proposition}[Boolean cube and reduced retract]
\label{prop:ds-boolean-retract}
The map $\Theta_m$ is an isomorphism of $A$-semimodules, where the scalar action
on $B_m$ is
\[
 \tau I=\varnothing,\qquad eI=I.
\]
The maps
\[
 \Sigma_m:A^m\to A,\quad \Sigma_m(a)=a_1+\cdots+a_m,
 \qquad
 \Delta_m:A\to A^m,\quad \Delta_m(a)=(a,\ldots,a),
\]
are $A$-linear and satisfy $\Sigma_m\Delta_m=\operatorname{id}_A$.  In subset
coordinates,
\[
 \Sigma_m(I)=
 \begin{cases}\tau,&I=\varnothing,\\ e,&I\ne\varnothing,\end{cases}
 \qquad
 \Delta_m(\tau)=\varnothing,\quad \Delta_m(e)=[m].
\]
Thus the only nontrivial fibre of $\Sigma_m$ is the set of all
$2^m-1$ nonempty subsets.
\end{proposition}

\begin{proof}
Coordinatewise addition in $A^m$ is Boolean OR, so
$\Theta_m(a+b)=\Theta_m(a)\cup\Theta_m(b)$; the inverse assigns $e$ to the
coordinates in a subset and $\tau$ to its complement.  The formulas for
$\Sigma_m$ and $\Delta_m$ are then the nonempty-set detector and the inclusion
of the two extreme vertices.  They preserve joins and the scalar action, and
their composite fixes both elements of $A$.  The fibre count is immediate from
the displayed formula.
\end{proof}

Applying these maps independently at each $x$ gives sheaf morphisms
\[
 \mathcal V_D\mathop{\longrightarrow}^{\Sigma_D}\mathcal A_D
 \mathop{\longrightarrow}^{\Delta_D}\mathcal V_D,
 \qquad
 \Sigma_D\Delta_D=\operatorname{id}_{\mathcal A_D}.
\]
If $\mathcal V_D\cong\mathcal A_D$, then their stalks at $x$ have respectively
$2^{m_x}$ and $2$ elements; hence $m_x=1$.  Conversely that condition makes
the definitions identical.  This proves, without erasing the retraction, that
\[
 \mathcal V_D\cong\mathcal A_D
 \quad\Longleftrightarrow\quad
 m_x=1\text{ for every }x\in Z.
\]

\subsection{The intrinsic Artin chain and both adjoints to the cube inclusion}

Let $\mathfrak m_x$ be the maximal ideal of the analytic local ring
$\mathcal O_{X,x}$ and define
\[
 Q_{x,m}:=\mathcal O_{X,x}/\mathfrak m_x^m.
\]
A centred holomorphic coordinate $t$ identifies
\[
 \mathcal O_{X,x}\cong\C\{t\},\qquad
 Q_{x,m}\cong\C\{t\}/(t^m)\cong\C[t]/(t^m).
\]

\begin{proposition}[Complete ideal lattice]
\label{prop:ds-artin-chain}
The ideals of $Q_{x,m}$ are exactly
\[
 I_k:=\mathfrak m_x^k/\mathfrak m_x^m,
 \qquad 0\leq k\leq m,
\]
and
\[
 Q_{x,m}=I_0\supset I_1\supset\cdots\supset I_m=0.
\]
This chain is intrinsic: changing the centred coordinate changes its polynomial
presentation but not any $I_k$.
\end{proposition}

\begin{proof}
Every nonzero convergent germ has a unique factorization $t^ku(t)$ with
$k\geq0$ and $u(0)\ne0$.  If $J\ne0$ is an ideal of $\C\{t\}$, choose the least
order $k$ of a nonzero element of $J$.  A least-order element is $t^ku$ with
$u$ a unit, so $t^k\in J$; every element of $J$ has order at least $k$, so
$J\subseteq(t^k)$.  Hence $J=(t^k)$.  Ideals of the quotient correspond to ideals
of $\C\{t\}$ containing $(t^m)$, which are precisely $(t^k)$ for
$0\leq k\leq m$.  Finally $(t)=\mathfrak m_x$, so the displayed ideals are the
coordinate-free powers of the maximal ideal.  This is the analytic instance of
the DVR ideal classification recorded in \cite[\S13.5, especially
Thm.~13.5.4]{vakil2024}.
\end{proof}

Define terminal segments
\[
 T_r:=\{r,r+1,\ldots,m\}\quad(1\leq r\leq m),
 \qquad T_{m+1}:=\varnothing,
\]
and let $C_m=\{T_1,\ldots,T_{m+1}\}\subset B_m$.  The intrinsic identification
\[
 \operatorname{Idl}(Q_{x,m})\longrightarrow C_m,
 \qquad I_k\longmapsto T_{k+1},
\]
preserves inclusion and joins.  Let $i_m:C_m\hookrightarrow B_m$ be inclusion,
and define
\[
 r_m(I)=
 \begin{cases}
 \varnothing,&I=\varnothing,\\
 T_{\min I},&I\ne\varnothing,
 \end{cases}
 \qquad
 s_m(I)=\{j\in[m]:T_j\subseteq I\}.
\]

\begin{theorem}[Coordinate-level chain--cube comparison]
\label{thm:ds-chain-cube}
For $I\in B_m$ and $T\in C_m$,
\[
 r_m(I)\subseteq T\iff I\subseteq i_m(T),
 \qquad
 i_m(T)\subseteq I\iff T\subseteq s_m(I).
\]
Thus $r_m\dashv i_m\dashv s_m$ in finite posets.  The map $r_m$ preserves
joins and $r_mi_m=\operatorname{id}_{C_m}$, so $C_m$ is a split join-subsemilattice
of $B_m$.  The fibres are exactly
\[
 r_m^{-1}(T_r)=
 \begin{cases}
 \{I\subseteq T_r:r\in I\},&1\leq r\leq m,\\
 \{\varnothing\},&r=m+1,
 \end{cases}
\]
and
\[
 s_m^{-1}(T_r)=
 \begin{cases}
 \{[m]\},&r=1,\\
 \{I:T_r\subseteq I,\ r-1\notin I\},&2\leq r\leq m+1.
 \end{cases}
\]
In the last line $T_{m+1}=\varnothing$, so the condition is $m\notin I$.
Consequently the two adjoints forget, respectively, every coordinate above the
first occupied coordinate and every coordinate below the terminal occupied run.
\end{theorem}

\begin{proof}
The set $r_m(I)$ is the least terminal segment containing $I$, while $s_m(I)$ is
the greatest terminal segment contained in $I$.  These two universal properties
are exactly the displayed adjunctions.  Also
\[
 r_m(I\cup J)=r_m(I)\cup r_m(J):
\]
if both sets are nonempty, both sides are
$T_{\min(\min I,\min J)}$, and the empty cases are immediate.  A terminal
segment is fixed by terminal closure, proving $r_mi_m=\operatorname{id}$.

For $1\leq r\leq m$, the equality $r_m(I)=T_r$ says precisely that $r$ is the
least member of $I$, equivalently $r\in I\subseteq T_r$.  The empty fibre is
obvious.  For $2\leq r\leq m$, the equality $s_m(I)=T_r$ says that the whole
tail $T_r$ lies in $I$ but the next longer tail does not; this is equivalent to
$T_r\subseteq I$ and $r-1\notin I$.  The endpoint cases follow from the same
description.  These formulas also exhibit the information loss rather than merely
counting it.
\end{proof}

For reference, the fibre cardinalities obtained from the proof are
\[
 |r_m^{-1}(T_r)|=2^{m-r}\ (r\leq m),
\]
and
\[
 |s_m^{-1}(T_1)|=1,\qquad
 |s_m^{-1}(T_r)|=2^{r-2}\ (2\leq r\leq m+1).
\]
When $m\geq2$, $C_m$ has width $1$ whereas $B_m$ contains the incomparable
vertices $\{1\}$ and $\{2\}$.  This is an explicit poset-isomorphism obstruction,
not an assertion that the objects have no relation: their positive relation is the
adjoint triple just proved.

Define product sheaves
\[
 \mathcal I_D(U):=\prod_{x\in U\cap Z}\operatorname{Idl}(Q_{x,m_x}),
 \qquad
 \mathcal B_D(U):=\prod_{x\in U\cap Z}B_{m_x}.
\]
The ideal-chain identification and $\Theta_{m_x}$ give
$\mathcal I_D\cong\prod C_{m_x}$ and identify $\mathcal B_D$ with
$\mathcal V_D$ after forgetting from $A$-semimodules to join-semilattices.  Applying
$i_{m_x}$ and $r_{m_x}$ coordinatewise gives morphisms of sheaves of
join-semilattices
\[
 \mathcal I_D\mathop{\longrightarrow}^{i_D}\mathcal B_D
 \mathop{\longrightarrow}^{r_D}\mathcal I_D,
 \qquad r_Di_D=\operatorname{id}_{\mathcal I_D},
\]
Applying $s_{m_x}$ coordinatewise gives a morphism of sheaves of posets
$s_D:\mathcal B_D\to\mathcal I_D$; it is not asserted to preserve joins.  Stalkwise
$r_D\dashv i_D\dashv s_D$.  Every map commutes with restriction
because restriction is coordinate projection.  This is the exact global morphism
between the intrinsic Artin model and the Boolean multiplicity model.

\subsection{Enhanced jets and the exact support morphism}

Now take $X=\C$ with its literal global coordinate $s$.  For $x\in Z$, put
$t_x=s-x$ and define
\[
 J_x:=\C[t_x]/(t_x^{m_x}),
 \qquad
 E_x:=t_x\C[t_x]/(t_x^{m_x+1}),
 \qquad
 L_x:=\C t_x^{m_x}\subset E_x.
\]
Thus $J_x$ has ordered basis $1,t_x,\ldots,t_x^{m_x-1}$ and $E_x$ has ordered
basis $t_x,\ldots,t_x^{m_x}$.  The coordinate $t_x$ gives an algebra isomorphism
$Q_{x,m_x}\cong J_x$ and a linear degree-shift isomorphism
\[
 d_x:E_x\longrightarrow J_x,
 \qquad
 d_x\left(\sum_{r=1}^{m_x}c_rt_x^r\right)
   =\sum_{r=1}^{m_x}c_rt_x^{r-1}.
\]
Its inverse is multiplication by $t_x$ followed by reduction modulo
$t_x^{m_x+1}$.  The map $d_x$ is a vector-space map, not a claim that $E_x$ has
the quotient-ring multiplication of $J_x$.

Define the coefficient-support map
\[
 \kappa_x:E_x\longrightarrow A^{m_x},
 \qquad
 \kappa_x\left(\sum_{r=1}^{m_x}c_rt_x^r\right)_r=
 \begin{cases}e,&c_r\ne0,\\ \tau,&c_r=0.\end{cases}
\]
and the basepoint-preserving section
\[
 \iota_x:A^{m_x}\longrightarrow E_x,
 \qquad
 \iota_x(a)=\sum_{\{r:a_r=e\}}t_x^r.
\]

\begin{proposition}[Support fibres and the additive obstruction]
\label{prop:ds-jet-support}
The maps $\kappa_x$ and $\iota_x$ are morphisms of pointed sets and
\[
 \kappa_x\iota_x=\operatorname{id}_{A^{m_x}}.
\]
For the vertex corresponding to $I\subseteq[m_x]$, the fibre is
\[
 \kappa_x^{-1}(I)=
 \left\{\sum_{r\in I}c_rt_x^r:c_r\in\C^\times\right\}
 \cong(\C^\times)^{|I|},
\]
with the empty fibre interpreted as the singleton $\{0\}$.  The map $\kappa_x$
is not additive.  Indeed, for every $m_x\geq1$,
\[
 \kappa_x(t_x)=\kappa_x(-t_x)=\{1\},
 \qquad
 \kappa_x(t_x-t_x)=\varnothing
 \ne\{1\}\cup\{1\}.
\]
Finally, the additive monoids $E_x$ and $A^{m_x}$ are not isomorphic: the latter
is idempotent, whereas $t_x+t_x=2t_x\ne t_x$ in characteristic zero.
\end{proposition}

\begin{proof}
The coefficient $1$ in $\iota_x(a)$ is nonzero in exactly the positions in which
$a_r=e$, proving the retraction identity.  Prescribing a support $I$ forces the
coefficients in $I$ to be nonzero and all other coefficients to vanish, which gives
the fibre formula.  The displayed cancellation proves nonadditivity.  An additive
monoid isomorphism preserves the sentence $v+v=v$ for every $v$; the final two
monoids satisfy different such sentences, yielding the stated obstruction.
\end{proof}

The product assignments
\[
 \mathcal J_D(U)=\prod_{x\in U\cap Z}J_x,
 \quad
 \mathcal E_D(U)=\prod_{x\in U\cap Z}E_x,
 \quad
 \mathcal L_D(U)=\prod_{x\in U\cap Z}L_x
\]
are sheaves by \cref{lem:ds-product-sheaf}.  The maps $d_x,\kappa_x,\iota_x$
assemble coordinatewise into sheaf maps, including the pointed-set retraction
\[
 \mathcal E_D\mathop{\longrightarrow}^{\kappa_D}\mathcal V_D
 \mathop{\longrightarrow}^{\iota_D}\mathcal E_D,
 \qquad \kappa_D\iota_D=\operatorname{id}_{\mathcal V_D}.
\]
Thus the additive obstruction and the positive pointed-set morphism are retained
simultaneously.

\begin{proposition}[Terminal jet and the Artin socle]
\label{prop:ds-terminal-socle}
Under $d_x:E_x\to J_x$, the line $L_x$ maps isomorphically to
\[
 \operatorname{Soc}(J_x)
 :=\{q\in J_x:t_xq=0\}
 =\C t_x^{m_x-1}.
\]
Its nonzero elements have Boolean support $T_{m_x}=\{m_x\}$, which corresponds
under \cref{thm:ds-chain-cube} to the smallest nonzero Artin ideal
$I_{m_x-1}$.
\end{proposition}

\begin{proof}
If $q=\sum_{r=0}^{m_x-1}c_rt_x^r$, then $t_xq=0$ modulo $t_x^{m_x}$ exactly
when every $c_r$ with $0\leq r\leq m_x-2$ vanishes; this condition is vacuous
when $m_x=1$.  Hence the socle is the displayed line.
The degree shift sends $t_x^{m_x}$ to $t_x^{m_x-1}$, while coefficient support
sends every nonzero scalar multiple of $t_x^{m_x}$ to $\{m_x\}=T_{m_x}$.
The ideal-chain coordinate sends $I_{m_x-1}$ to the same terminal segment.
\end{proof}

\begin{theorem}[Local leading class with the required hypothesis]
\label{thm:ds-leading-class}
Let $f\not\equiv0$ be meromorphic on $\C$, let $Z_f$ be its zero set, and assume only that
$f$ is holomorphic in a neighbourhood of every $x\in Z_f$.  If
$m_x=\operatorname{ord}_x(f)$, then its truncated class in $E_x$ is
\[
 [f]_x=\frac{f^{(m_x)}(x)}{m_x!}\,t_x^{m_x}
 \in L_x\setminus\{0\}.
\]
Its coefficient support is $T_{m_x}$, and its degree shift is a nonzero generator
of the Artin socle.
\end{theorem}

\begin{proof}
Holomorphy near $x$ and the definition of vanishing order give the Taylor
expansion
\[
 f(s)=\sum_{r\geq m_x}\frac{f^{(r)}(x)}{r!}t_x^r,
 \qquad f^{(m_x)}(x)\ne0.
\]
Reduction modulo $t_x^{m_x+1}$ leaves precisely the displayed term.  The support
and socle claims are \cref{prop:ds-terminal-socle}.  No global entire-function
hypothesis is used.
\end{proof}

\subsection{Compactification and the tail stalk}

Let $j:\C\hookrightarrow\mathbb P^1(\C)$ be the standard inclusion and let
$i_\infty:\{\infty\}\hookrightarrow\mathbb P^1(\C)$.  For a family of pointed
sets $(M_x)_{x\in Z}$, define
\[
 \operatorname{Tail}(M):=\left(\prod_{x\in Z}M_x\right)\!\big/\sim,
 \qquad
 a\sim b\iff a_x=b_x\text{ for all but finitely many }x.
\]
For vector spaces $V_x$, this is the vector quotient
\[
 \operatorname{Tail}_{\C}(V)
 =\left(\prod_{x\in Z}V_x\right)\Big/\left(\bigoplus_{x\in Z}V_x\right).
\]
Let $j_*\mathcal F_M(U)=\mathcal F_M(U\cap\C)$.  Define explicitly
$j^{\mathrm{fin}}_!\mathcal F_M$ by the same formula when $\infty\notin U$ and,
when $\infty\in U$, by retaining exactly the families with finite support.

\begin{lemma}[The finite extension is a sheaf]
\label{lem:ds-jfin-sheaf}
The assignment $j^{\mathrm{fin}}_!\mathcal F_M$ is a sheaf, and its inclusion into
$j_*\mathcal F_M$ is an isomorphism away from $\infty$.
\end{lemma}

\begin{proof}
Only gluing over an open set $U$ containing $\infty$ requires proof.  Let
$U=\bigcup_iU_i$ and let finite-support sections on the $U_i$ be compatible.
Choose $i_0$ with $\infty\in U_{i_0}$.  The complement
$K=\mathbb P^1(\C)\setminus U_{i_0}$ is a compact subset of $\C$.  Because $Z$
is closed discrete, $K\cap Z$ is finite.  The section on $U_{i_0}$ has finite
support, and compatibility forces the glued product section to vanish at every
point of $(U\cap Z)\cap U_{i_0}$ outside that finite support.  The only remaining
possible support lies in $K\cap Z$, also finite.  Hence the glued section has
finite support.  Locality and the statement away from $\infty$ follow directly
from the definition.
\end{proof}

\begin{theorem}[Exact tail stalk]
\label{thm:ds-tail-stalk}
The stalks of $j^{\mathrm{fin}}_!\mathcal F_M\to j_*\mathcal F_M$ agree at every
finite point.  At infinity,
\[
 (j^{\mathrm{fin}}_!\mathcal F_M)_\infty\cong *,
 \qquad
 (j_*\mathcal F_M)_\infty\cong\operatorname{Tail}(M).
\]
For vector spaces there is a short exact sequence
\[
 0\longrightarrow j^{\mathrm{fin}}_!\mathcal F_V
 \longrightarrow j_*\mathcal F_V
 \longrightarrow i_{\infty*}\operatorname{Tail}_{\C}(V)
 \longrightarrow0.
\]
Both the middle and right sheaves are flasque.  Therefore
\[
 H^q\!\left(\mathbb P^1(\C),j^{\mathrm{fin}}_!\mathcal F_V\right)=0
 \quad(q\geq1),
\]
and
\[
 \operatorname{Tail}_{\C}(V)
 \cong
 \Gamma(j_*\mathcal F_V)\big/\Gamma(j^{\mathrm{fin}}_!\mathcal F_V).
\]
\end{theorem}

\begin{proof}
The finite stalks are computed as in \cref{lem:ds-product-sheaf}.  A finite-support
germ near infinity becomes the basepoint after shrinking past its finite support,
so the first infinite stalk is terminal.  A neighbourhood of infinity has compact
complement in $\C$, and therefore meets $Z$ in a cofinite subset.  Extending a
family by basepoints across the finite complement identifies a germ of
$j_*\mathcal F_M$ with a global family modulo eventual equality.  This proves the
stalk formula.

For vector spaces, the quotient map sends a family on a neighbourhood of
$\infty$ to its class after zero extension over the finite complement.  It is
surjective, and its kernel is exactly finite support, proving the sheaf sequence.
Every restriction map of $j_*\mathcal F_V$ is a coordinate projection and is
surjective by zero extension.  The same is immediate for the skyscraper at
$\infty$; hence both are flasque.  The flasque-acyclicity theorem
\cite[pp.~184--188, Thm.~6.34]{lee2024} and the long exact sequence give
vanishing in degrees at least two.  In degree one, the relevant segment is
\[
 \Gamma(j_*\mathcal F_V)\longrightarrow
 \operatorname{Tail}_{\C}(V)\longrightarrow
 H^1(j^{\mathrm{fin}}_!\mathcal F_V)\longrightarrow0.
\]
The first arrow is the defining product-to-product/direct-sum quotient and is
surjective.  Thus $H^1$ vanishes as well, and its kernel gives the global-section
quotient formula.
\end{proof}

\subsection{The zeta divisor, its jet section, and its nonzero tails}

The Riemann zeta function is meromorphic on $\C$, with its only pole a simple
pole at $1$ \cite[\S25.2(i)]{apostolDLMF25}.  Define
\[
 Z_\zeta=\{\rho\in\C:\zeta(\rho)=0\},
 \qquad
 D_\zeta=\sum_{\rho\in Z_\zeta}m_\rho[\rho],
 \quad m_\rho=\operatorname{ord}_\rho(\zeta).
\]
The zero set is closed and discrete: every zero has a neighbourhood on which
$\zeta$ is holomorphic and not identically zero, so local factorization makes the
zero isolated; the pole at $1$ is not a finite accumulation point of zeros.  Thus
all preceding sheaves apply to $D_\zeta$.

\begin{theorem}[Zeta divisor shadow and global leading jet]
\label{thm:ds-zeta-shadow-jet}
At every $\rho\in Z_\zeta$,
\[
 (\mathcal V_{D_\zeta})_\rho\cong A^{m_\rho},
 \qquad
 (\mathcal A_{D_\zeta})_\rho\cong A,
\]
and the second stalk is the literal fixed set from
\cref{prop:ds-fixed-pair}.  The family
\[
 [\zeta]_\rho
 =\lambda_\rho(s-\rho)^{m_\rho},
 \qquad
 \lambda_\rho=\frac{\zeta^{(m_\rho)}(\rho)}{m_\rho!}\ne0,
\]
is a global section of $\mathcal E_{D_\zeta}$.  Its support at $\rho$ is
$T_{m_\rho}$ and its degree shift generates the socle of
$\C[s-\rho]/(s-\rho)^{m_\rho}$.
\end{theorem}

\begin{proof}
The stalk formulas are \cref{lem:ds-product-sheaf}; the fixed-set statement is
\cref{prop:ds-fixed-pair}.  Although $\zeta$ is not entire, it is holomorphic near
each of its zeros.  Hence \cref{thm:ds-leading-class} applies at every $\rho$.
Arbitrary families are sections of the product sheaf, which proves globality.
The support and socle conclusions were proved in
\cref{prop:ds-terminal-socle}.
\end{proof}

The reflection formula
\[
 \zeta(s)=2(2\pi)^{s-1}\sin\!\left(\frac{\pi s}{2}\right)
 \Gamma(1-s)\zeta(1-s)
\]
is DLMF equation 25.4.2 \cite[\S25.4, Eq.~25.4.2]{apostolDLMF25}; it also
implies the trivial zeros \cite[\S25.10(i)]{apostolDLMF25}.

\begin{proposition}[Trivial-zero multiplicity and leading coordinate]
\label{prop:ds-trivial-derivative}
For every $n\geq1$, the zero $-2n$ is simple and
\[
 \zeta'(-2n)
 =(-1)^n\frac{(2n)!}{2(2\pi)^{2n}}\,\zeta(2n+1).
\]
\end{proposition}

\begin{proof}
At $s=-2n$, every factor in the reflection formula except the sine is holomorphic
and nonzero: in particular $\Gamma(1+2n)=(2n)!$ and
$\zeta(1+2n)>0$ by its positive Dirichlet series.  The sine has derivative
$(\pi/2)\cos(-n\pi)=(\pi/2)(-1)^n$, so the zero is simple.  Differentiating the
product at that simple sine zero leaves
\[
 2(2\pi)^{-2n-1}\Gamma(1+2n)\zeta(1+2n)
 \cdot\frac\pi2(-1)^n,
\]
which is the displayed expression after collecting powers of $2$ and $\pi$.
\end{proof}

\begin{corollary}[Nontriviality of both tail objects]
\label{cor:ds-zeta-tail}
The Boolean tail
\[
 \operatorname{Tail}\bigl((A^{m_\rho})_{\rho\in Z_\zeta}\bigr)
\]
has more than one element, and the vector tail
\[
 \left(\prod_{\rho\in Z_\zeta}E_\rho\right)
 \Big/\left(\bigoplus_{\rho\in Z_\zeta}E_\rho\right)
\]
contains the nonzero class of $([\zeta]_\rho)_\rho$.
\end{corollary}

\begin{proof}
The trivial zeros $-2,-4,\ldots$ form an infinite subset.  A family equal to a
fixed nonbasepoint at every trivial-zero coordinate is not eventually the
basepoint, so it gives a nonbasepoint Boolean tail class.  By
\cref{prop:ds-trivial-derivative}, every coordinate $[\zeta]_{-2n}$ is nonzero.
Hence the jet family is not finitely supported and its class in the vector quotient
is nonzero.
\end{proof}

\subsection{Equivariance over nonidentity base maps}

Before specializing to the completed function, we state the base-map convention
explicitly.  Let $h:X\to X$ be a homeomorphism satisfying $h(Z)=Z$ and
$m_{h(x)}=m_x$ for every $x\in Z$.  For every open $U\subset X$, define
\[
 T_h(U):\mathcal V_D(U)\longrightarrow\mathcal V_D(hU),
 \qquad
 \bigl(T_h(U)a\bigr)_{h(x),r}=a_{x,r}
 \quad(1\leq r\leq m_x).
\]
The same coordinate transport defines $T_h$ on $\mathcal I_D$.

\begin{proposition}[Exact transport and cocycle]
\label{prop:ds-equivariant-transport}
Every $T_h(U)$ is an isomorphism, with inverse $T_{h^{-1}}(hU)$, and
\[
 \operatorname{res}_{hU,hV}T_h(U)
 =T_h(V)\operatorname{res}_{U,V}
 \quad(V\subset U).
\]
For two divisor-preserving homeomorphisms $g,h$,
\[
 T_g(hU)T_h(U)=T_{g\circ h}(U).
\]
Thus these maps are equivariant sheaf data over the displayed base maps.  They
are not being silently recast as endomorphisms over $\operatorname{id}_X$.
Moreover the Artin--Boolean comparison is equivariant:
\[
 T_h i_D=i_D T_h,\qquad T_h r_D=r_D T_h,\qquad T_h s_D=s_D T_h.
\]
\end{proposition}

\begin{proof}
Each formula merely relabels the coordinate $(x,r)$ as $(h(x),r)$, so the inverse
and restriction identities follow by substitution.  Relabelling first by $h$ and
then by $g$ relabels by $g\circ h$, proving the cocycle.  The three chain--cube
maps act on the $r$-coordinates at one support point and do not change the support
point; relabelling and applying them therefore commute coordinate by coordinate.
\end{proof}

\subsection{Completed-zeta packet descent and phase coordinates}

Use Riemann's completed function in the literal convention
\[
 \xi(s)=\frac12s(s-1)\pi^{-s/2}\Gamma\!\left(\frac s2\right)\zeta(s).
\]
DLMF gives this normalization and the symmetry $\xi(1-s)=\xi(s)$
\cite[\S25.4, Eqs.~25.4.3--25.4.4]{apostolDLMF25}.  Lagarias and Montague
record explicitly that $\xi$ is entire and that its zeros are exactly the zeta
zeros in $0<\Re(s)<1$
\cite[\S1, Eq.~(1.1) and the following paragraph]{lagariasMontague2011}.
On this open strip the remaining factor
$\tfrac12s(s-1)\pi^{-s/2}\Gamma(s/2)$ is holomorphic and nonzero, so the
vanishing orders of $\xi$ and $\zeta$ agree.
Let $D_\xi$ be this divisor and set
\[
 c(s)=\overline s,\qquad w(s)=1-s,\qquad
 W=\{1,c,w,cw\}\cong C_2\times C_2.
\]
The identity $\overline{\xi(\overline s)}=\xi(s)$ follows first in a right
half-plane from the real Dirichlet coefficients and the conjugation law for
$\Gamma$, and then everywhere by analytic continuation.  Hence both $c$ and
$w$ preserve $D_\xi$, including multiplicities.

Let $Y_\xi=|D_\xi|/W$ with quotient map $q:|D_\xi|\to Y_\xi$.  Since the source
is discrete and $W$ is finite, $Y_\xi$ is discrete.  For an orbit $O$, write
$r_O=|O|$ and $m_O$ for its common multiplicity.

\begin{theorem}[Packet descent with explicit diagonal inverse]
\label{thm:ds-packet-descent}
On the restriction of the divisor sheaves to $|D_\xi|$,
\[
 (q_*\mathcal V_{D_\xi})_O\cong(A^{m_O})^{r_O},
 \qquad
 (q_*\mathcal I_{D_\xi})_O\cong(C_{m_O})^{r_O}.
\]
The fixed parts under the transport action of $W$ are the diagonal subobjects.
Projection to any chosen factor and diagonal repetition are mutually inverse on
the fixed part.  Consequently
\[
 (q_*\mathcal V_{D_\xi})^W(U)
 \cong\prod_{O\in U}A^{m_O},
 \qquad
 (q_*\mathcal I_{D_\xi})^W(U)
 \cong\prod_{O\in U}C_{m_O}
\]
for every open $U\subseteq Y_\xi$.
The descended chain--cube maps still satisfy
$r\dashv i\dashv s$ stalkwise and $ri=\operatorname{id}$.
\end{theorem}

\begin{proof}
Because the orbit space is discrete, the stalk of a pushforward at $O$ is the
product of the stalks over $q^{-1}(O)=O$.  Multiplicity preservation gives the
two displayed products.  The $W$-action permutes the factors transitively.
A tuple is fixed exactly when every factor equals every other factor, which is
the diagonal.  Projection of a diagonal tuple to one component is independent of
the chosen component, and diagonal repetition is its two-sided inverse.  These
maps commute with restrictions, so the stalkwise descriptions assemble into the
displayed product sheaves.  The equivariance equations in
\cref{prop:ds-equivariant-transport} descend the three comparison maps and their
equations.
\end{proof}

For a zero $\rho$ of $\xi$, define the nonzero leading coefficient and its phase
\[
 \mu_\rho:=\frac{\xi^{(m_\rho)}(\rho)}{m_\rho!}\in\C^\times,
 \qquad
 \phi_\rho:=\frac{\mu_\rho}{|\mu_\rho|}\in S^1.
\]

\begin{proposition}[Exact phase transport]
\label{prop:ds-phase-transport}
For every such $\rho$,
\[
 \mu_{\overline\rho}=\overline{\mu_\rho},
 \qquad
 \mu_{1-\rho}=(-1)^{m_\rho}\mu_\rho,
 \qquad
 \mu_{1-\overline\rho}=(-1)^{m_\rho}\overline{\mu_\rho},
\]
and the identical formulas hold for $\phi$.
\end{proposition}

\begin{proof}
Expand $\overline{\xi(\overline s)}=\xi(s)$ in local power series at $\rho$;
coefficient conjugation gives the first formula.  Differentiating
$\xi(1-s)=\xi(s)$ exactly $m_\rho$ times gives
\[
 (-1)^{m_\rho}\xi^{(m_\rho)}(1-s)=\xi^{(m_\rho)}(s).
\]
Evaluation at $s=\rho$ and division by $m_\rho!$ gives the second formula; the
third is its conjugate.  Absolute values are preserved by conjugation and by the
sign, so division by absolute value gives the phase formulas.
\end{proof}

For each orbit $O$, let $\Gamma_{m_O}$ be the subgroup of
$\operatorname{Homeo}(S^1)$ generated by
\[
 z\longmapsto\overline z,\qquad z\longmapsto(-1)^{m_O}z,
\]
and put $\Phi_O=S^1/\Gamma_{m_O}$.  By
\cref{prop:ds-phase-transport}, the image of $\phi_\rho$ in $\Phi_O$ is independent
of $\rho\in O$.  These images form a section of the product sheaf
\[
 \mathcal{Ph}_\xi(U)=\prod_{O\in U}\Phi_O
\]
on $Y_\xi$.  This is a quotient-valued phase decoration of the packet sheaf; no
claim about the location, multiplicity, or spectral origin of any individual zero
is inferred from it.

\subsection{Certificate boundary and retained interfaces}

The accompanying exact finite certificate exhaustively checks, for
$1\leq m\leq8$, the two adjunctions, the retraction, join preservation, every
fibre formula and fibre cardinality, the reduced-shadow retraction, coefficient
support fibres over a finite coefficient alphabet, and the Klein-four phase-action
relations.  Those checks are a replay of finite algebra, not a replacement for the
all-$m$ proofs above.  In particular, they do not certify analytic continuation,
the functional equation, divisor discreteness, sheaf cohomology, or any statement
about individual zeros.

Three programme boundaries remain explicit.  First, $A^{m_x}$ is a constructed
multiplicity shadow, while $Q_{x,m_x}$ is the intrinsic analytic thickening; their
proved interface is $r_m\dashv i_m\dashv s_m$ and not an untyped identification.
Second, the coefficient-support map from jets is a pointed-set retraction with
explicit torus fibres and a proved additive obstruction.  Third, the symmetry maps
are transport over $c$ and $w$, followed by fixed-point descent along $q$; they are
not endomorphisms over an unchanged base.  Later blueprint, analytic packet-product,
operator-algebraic, or physical interpretations require their own typed arrows and
are not supplied by the present construction.
